 >>?\chapter{Revis o de Geometria Diferencial}
\label{apx:geometria}

Este ap andice apresenta os conceitos fundamentais de geometria diferencial necess irios compreens o da Abordagem de Apre SSamento por Transporte Geom ctrico (GAT). O leitor familiarizado com a linguagem de fibrados e formas diferenciais pode consult i-lo como refer ancia pontual; o leitor iniciante encontrar i aqui as defini SS ues essenciais, enunciados dos principais teoremas e nota SS o empregada ao longo do texto. Seguimos de perto a exposi SS o cl issica de \citeonline{nakahara2003geometry}, complementada por \citeonline{frankel2011geometry} e \citeonline{bleecker1981gauge}.

%%%%%%%%%%%%%%%%%%%%%%%%%%%%%%%%%%%%%%%%%%%%%%%%%%%%%%%%%%%%%%%%%%%%%%%%%%%%%%%%
\section{Variedades Diferenci iveis}
\label{sec:variedades}

\subsection{Cartas e Atlas}

\begin{defi}[Carta Local]
Seja $M$ um espa SSo topol 3gico Hausdorff com base enumer ivel. Uma \textbf{carta local} (ou sistema de coordenadas) em $M$ c um par $(U,\varphi)$, em que $U\subset M$ c um aberto e $\varphi:U\to\varphi(U)\subset\R^n$ c um homeomorfismo sobre sua imagem. As fun SS ues $x^i = u^i\circ\varphi:U\to\R$, em que $u^i:\R^n\to\R$ s o as proje SS ues can 'nicas, denominam-se \textbf{coordenadas locais}.
\end{defi}

\begin{defi}[Atlas]
Um \textbf{atlas} $\mathcal{A}$ de classe $C^k$ sobre $M$ c uma cole SS o de cartas $\{(U_\alpha,\varphi_\alpha)\}_{\alpha\in I}$ tal que:
\begin{enumerate}
 \item $\bigcup_{\alpha\in I}U_\alpha = M$ (cobertura);
 \item Para todo $\alpha,\beta\in I$ com $U_\alpha\cap U_\beta\neq\varnothing$, a \textbf{mudan SSa de coordenadas}
 \begin{equation}
 \varphi_\beta\circ\varphi_\alpha^{-1}:\varphi_\alpha(U_\alpha\cap U_\beta)\to\varphi_\beta(U_\alpha\cap U_\beta)
 \end{equation}
 c um difeomorfismo de classe $C^k$ entre abertos de $\R^n$.
\end{enumerate}
Duas cartas s o \textbf{compat -veis} se a uni o de seus atlas ainda constitui um atlas $C^k$. Um atlas c \textbf{m iximo} quando cont cm todas as cartas compat -veis com suas cartas.
\end{defi}

\begin{defi}[Variedade Diferenci ivel]
Uma \textbf{variedade diferenci ivel} de dimens o $n$ e classe $C^k$ c um par $(M,\mathcal{A})$, em que $M$ c um espa SSo topol 3gico Hausdorff com base enumer ivel e $\mathcal{A}$ c um atlas m iximo $C^k$ sobre $M$. Quando $k=\infty$, dizemos que $M$ c uma variedade \textbf{suave}.
\end{defi}

\begin{ex}
O espa SSo euclidiano $\R^n$ com a carta identidade $(\R^n,\id)$ c uma variedade suave. A esfera $S^n = \{x\in\R^{n+1}:\|x\|=1\}$ admite um atlas com duas cartas estereogr ificas, constituindo uma variedade suave compacta de dimens o $n$.
\end{ex}

\subsection{Espa SSo Tangente}

\begin{defi}[Vetor Tangente via Deriva SS ues]
Seja $p\in M$. Uma \textbf{deriva SS o} em $p$ c um funcional linear $X_p:C^\infty(M)\to\R$ satisfazendo a regra de Leibniz:
\begin{equation}
 X_p(fg) = X_p(f)\,g(p) + f(p)\,X_p(g),\qquad \forall f,g\in C^\infty(M).
\end{equation}
O \textbf{espa SSo tangente} $T_pM$ c o conjunto de todas as deriva SS ues em $p$, dotado da estrutura natural de espa SSo vetorial real de dimens o $n$.
\end{defi}

Em coordenadas locais $(U,\varphi)$ com $\varphi=(x^1,\dots,x^n)$, os operadores
\begin{equation}
 \left.\frac{\partial}{\partial x^i}\right|_p(f) = \frac{\partial (f\circ\varphi^{-1})}{\partial u^i}\Big|_{\varphi(p)}
\end{equation}
formam uma base de $T_pM$, denominada \textbf{base coordenada}. Todo vetor tangente expressa-se unicamente como $X_p = X^i\,\partial/\partial x^i|_p$ (soma de Einstein impl -cita).

\subsection{Campos Vetoriais e Colchete de Lie}

\begin{defi}
Um \textbf{campo vetorial} $X$ sobre $M$ c uma se SS o suave do fibrado tangente $TM$, isto c, uma aplica SS o $X:M\to TM$ tal que $\pi\circ X = \id_M$. Em coordenadas locais, $X = X^i(x)\,\partial/\partial x^i$.
\end{defi}

O conjunto $\mathfrak{X}(M)$ de todos os campos vetoriais suaves constitui um m 3dulo sobre $C^\infty(M)$ e uma ilgebra de Lie real com o \textbf{colchete de Lie}:
\begin{equation}
 [X,Y](f) = X(Y(f)) - Y(X(f)),\qquad f\in C^\infty(M).
\end{equation}
Em coordenadas, $[X,Y] = (X^j\partial_jY^i - Y^j\partial_jX^i)\,\partial_i$.

O colchete de Lie satisfaz a \textbf{identidade de Jacobi}:
\begin{equation}
 [X,[Y,Z]] + [Y,[Z,X]] + [Z,[X,Y]] = 0,
\end{equation}
e constitui o prot 3tipo alg cbrico da curvatura em geometria diferencial.

%%%%%%%%%%%%%%%%%%%%%%%%%%%%%%%%%%%%%%%%%%%%%%%%%%%%%%%%%%%%%%%%%%%%%%%%%%%%%%%%
\section{Formas Diferenciais}
\label{sec:formas}

\subsection{\texorpdfstring{$1$}{1}-Formas e o Fibrado Cotangente}

\begin{defi}
O \textbf{espa SSo cotangente} $T_p^*M$ c o dual do espa SSo tangente: $T_p^*M = (T_pM)^*$. Uma \textbf{1-forma diferencial} $\omega$ c uma se SS o suave do fibrado cotangente $T^*M$, atribuindo a cada $p\in M$ um funcional linear $\omega_p:T_pM\to\R$.
\end{defi}

A \textbf{diferencial exterior} de uma fun SS o $f\in C^\infty(M)$ c a 1-forma $df$ definida por $df(X) = X(f)$. Em coordenadas, $df = (\partial_i f)\,dx^i$, em que $\{dx^i\}$ c a base dual de $\{\partial/\partial x^i\}$, satisfazendo $dx^i(\partial_j) = \delta^i_j$.

\subsection{\texorpdfstring{$k$}{k}-Formas e o Produto Exterior}

Seja $\bigwedge^k T_p^*M$ a $k$- csima pot ancia exterior do espa SSo cotangente. Uma \textbf{$k$-forma diferencial} $\omega$ c uma se SS o do fibrado $\bigwedge^k T^*M$. Localmente,
\begin{equation}
 \omega = \frac{1}{k!}\,\omega_{i_1\cdots i_k}\,dx^{i_1}\wedge\cdots\wedge dx^{i_k},
\end{equation}
com coeficientes $\omega_{i_1\cdots i_k}$ totalmente antissim ctricos.

O \textbf{produto exterior} $\wedge:\Omega^k(M)\times\Omega^\ell(M)\to\Omega^{k+\ell}(M)$ c bilinear, associativo, e satisfaz $\alpha\wedge\beta = (-1)^{k\ell}\beta\wedge\alpha$ para $\alpha\in\Omega^k(M)$, $\beta\in\Omega^\ell(M)$.

\subsection{Derivada Exterior}

\begin{defi}
A \textbf{derivada exterior} $d:\Omega^k(M)\to\Omega^{k+1}(M)$ c a onica aplica SS o $\R$-linear satisfazendo:
\begin{enumerate}
 \item $df$ coincide com a diferencial de fun SS ues para $k=0$;
 \item $d(d\omega)=0$ para toda forma $\omega$ (nilpot ancia: $d^2=0$);
 \item $d(\alpha\wedge\beta) = d\alpha\wedge\beta + (-1)^k\alpha\wedge d\beta$ (regra de Leibniz graduada).
\end{enumerate}
Em coordenadas, para $\omega = \frac{1}{k!}\omega_{i_1\cdots i_k}dx^{i_1}\wedge\cdots\wedge dx^{i_k}$:
\begin{equation}
 d\omega = \frac{1}{k!}\,\partial_j\omega_{i_1\cdots i_k}\,dx^j\wedge dx^{i_1}\wedge\cdots\wedge dx^{i_k}.
\end{equation}
\end{defi}

\begin{defi}
Uma forma $\omega$ c \textbf{fechada} se $d\omega=0$, e \textbf{exata} se existe $\eta$ tal que $\omega = d\eta$. Toda forma exata c fechada, pois $d^2=0$.
\end{defi}

\subsection{Lema de Poincar c}

\begin{teo}[Lema de Poincar c]
Seja $U\subset\R^n$ um aberto contr itil (em particular, estrelado). Toda $k$-forma fechada em $U$ ($k\geq 1$) c exata. Equivalentemente, os grupos de cohomologia de de Rham de $U$ satisfazem $H^k_{\text{dR}}(U)=0$ para $k\geq 1$.
\end{teo}

\begin{proof}[Esbo SSo]
Define-se o operador de homotopia $K:\Omega^k(U)\to\Omega^{k-1}(U)$ por integra SS o ao longo de raios:
\begin{equation}
 (K\omega)_{i_1\cdots i_{k-1}}(x) = \int_0^1 t^{k-1}\,\omega_{j i_1\cdots i_{k-1}}(tx)\,x^j\,dt,
\end{equation}
que satisfaz $Kd + dK = \id$ em $\Omega^k(U)$ para $k\geq 1$. Se $d\omega=0$, ent o $\omega = d(K\omega)$, donde $\omega$ c exata.
\end{proof}

O Lema de Poincar c c fundamental na GAT: ele garante que toda 2-forma de curvatura fechada (identidade de Bianchi) c localmente exata, fornecendo potenciais de conex o (1-formas de gauge) em vizinhan SSas contr iteis do espa SSo de ativos.

%%%%%%%%%%%%%%%%%%%%%%%%%%%%%%%%%%%%%%%%%%%%%%%%%%%%%%%%%%%%%%%%%%%%%%%%%%%%%%%%
\section{Grupos de Lie}
\label{sec:grupos_lie}

\subsection{Defini SS o e Exemplos}

\begin{defi}
Um \textbf{grupo de Lie} $G$ c uma variedade suave dotada de uma estrutura de grupo tal que as opera SS ues de multiplica SS o $\mu:G\times G\to G$, $\mu(g,h)=gh$, e invers o $\iota:G\to G$, $\iota(g)=g^{-1}$, s o aplica SS ues suaves.
\end{defi}

\begin{ex}
O grupo linear geral $\GL(n,\R)$ c um grupo de Lie de dimens o $n^2$. Seus subgrupos fechados --- $\SL(n,\R)$, $\SO(n)$, $\U(n)$, $\SU(n)$ --- s o igualmente grupos de Lie (teorema de Cartan--von Neumann). Em finan SSas, $\GL(n,\R)$ modela transforma SS ues lineares invert -veis entre cestas de ativos, enquanto $\SO(n)$ preserva a estrutura de correla SS o (m ctrica de Mahalanobis).
\end{ex}

\subsection{ lgebra de Lie}

\begin{defi}
A \textbf{ ilgebra de Lie} $\mathfrak{g}$ de um grupo de Lie $G$ c o espa SSo tangente na identidade, $\mathfrak{g}=T_eG$, munido do colchete de Lie induzido pelos campos invariantes esquerda.
\end{defi}

Um campo $X\in\mathfrak{X}(G)$ c \textbf{invariante esquerda} se $(L_g)_*X = X$ para todo $g\in G$, em que $L_g(h)=gh$. O colchete de dois campos invariantes esquerda c invariante esquerda, conferindo a $\mathfrak{g}$ a estrutura de ilgebra de Lie.

Em termos matriciais, para $G\subset\GL(n,\R)$, tem-se $\mathfrak{g}\subset\mathfrak{gl}(n,\R)=M_n(\R)$ e $[A,B]=AB-BA$ (comutador de matrizes).

\subsection{Aplica SS o Exponencial}

\begin{defi}
A \textbf{aplica SS o exponencial} $\exp:\mathfrak{g}\to G$ c definida por $\exp(A)=\gamma_A(1)$, em que $\gamma_A:\R\to G$ c o subgrupo a 1-par metro gerado por $A\in\mathfrak{g}$, i.e., a onica curva satisfazendo $\gamma_A(0)=e$ e $\dot\gamma_A(t) = (L_{\gamma_A(t)})_*A$.
\end{defi}

Para grupos matriciais, $\exp(A) = \sum_{k=0}^\infty A^k/k!$, a exponencial de matrizes usual. A exponencial c um difeomorfismo local entre uma vizinhan SSa de $0\in\mathfrak{g}$ e uma vizinhan SSa de $e\in G$, fornecendo a \textbf{carta exponencial} fundamental na GAT para parametrizar instrumentos financeiros como elementos de grupos de Lie.

%%%%%%%%%%%%%%%%%%%%%%%%%%%%%%%%%%%%%%%%%%%%%%%%%%%%%%%%%%%%%%%%%%%%%%%%%%%%%%%%
\section{Fibrados: Defini SS ues, Conex ues e Curvatura}
\label{sec:fibrados}

\subsection{Defini SS o e Se SS ues}

\begin{defi}
Um \textbf{fibrado} (ou espa SSo fibrado) $\pi:E\to M$ consiste em um espa SSo total $E$, uma base $M$ (variedade), e uma proje SS o $\pi$ sobrejetora e suave. Para cada $x\in M$, a \textbf{fibra} $E_x=\pi^{-1}(x)$ c uma variedade difeomorfa a uma \textbf{fibra t -pica} $F$. Localmente, $E$ c trivial: existe uma vizinhan SSa $U\subset M$ e um difeomorfismo $\phi:\pi^{-1}(U)\to U\times F$ tal que $\pr_1\circ\phi = \pi$.
\end{defi}

\begin{ex}[Fibrados Relevantes na GAT]
\begin{itemize}
 \item \textbf{Fibrado tangente} $TM$: fibra t -pica $\R^n$, base $M$.
 \item \textbf{Fibrado cotangente} $T^*M$: dual de $TM$.
 \item \textbf{Fibrado de frames} $FM$: fibra em $x$ c o conjunto de todas as bases de $T_xM$; isomorfo a $\GL(n,\R)$.
 \item \textbf{Fibrado principal} $P(M,G)$: fibra t -pica c um grupo de Lie $G$ agindo livremente direita.
 \item \textbf{Fibrado associado} $E = P\times_G V$: obtido de um fibrado principal $P$ e uma representa SS o $\rho:G\to\GL(V)$.
\end{itemize}
\end{ex}

Uma \textbf{se SS o} de $\pi:E\to M$ c uma aplica SS o suave $s:M\to E$ tal que $\pi\circ s = \id_M$. Campos vetoriais s o se SS ues de $TM$; 1-formas s o se SS ues de $T^*M$.

\subsection{Conex ues e Derivada Covariante}

\begin{defi}
Uma \textbf{conex o} (ou derivada covariante) em um fibrado vetorial $E\to M$ c um operador $\nabla:\mathfrak{X}(M)\times\Gamma(E)\to\Gamma(E)$, $(X,s)\mapsto\nabla_X s$, $\R$-bilinear e satisfazendo:
\begin{enumerate}
 \item $\nabla_{fX}s = f\nabla_X s$ ($C^\infty(M)$-linearidade no argumento do campo);
 \item $\nabla_X(fs) = X(f)s + f\nabla_X s$ (regra de Leibniz no argumento da se SS o).
\end{enumerate}
\end{defi}

Em uma base local de se SS ues $\{e_a\}$ de $E$ sobre $U\subset M$, a conex o c determinada pela \textbf{1-forma de conex o} $\omega$:
\begin{equation}
 \nabla_X e_a = \omega^b_{\;a}(X)\,e_b,
\end{equation}
em que $\omega = (\omega^b_{\;a})$ c uma matriz de 1-formas diferenciais em $U$, denominada forma de conex o local.

\subsection{Curvatura}

\begin{defi}
A \textbf{curvatura} de uma conex o $\nabla$ c o operador $R:\mathfrak{X}(M)\times\mathfrak{X}(M)\times\Gamma(E)\to\Gamma(E)$ definido por:
\begin{equation}
 R(X,Y)s = \nabla_X\nabla_Y s - \nabla_Y\nabla_X s - \nabla_{[X,Y]}s.
\end{equation}
\end{defi}

Em termos da forma de conex o, a curvatura c representada pela \textbf{2-forma de curvatura} $\Omega$:
\begin{equation}
 \Omega^a_{\;b} = d\omega^a_{\;b} + \omega^a_{\;c}\wedge\omega^c_{\;b}.
\end{equation}
Esta c a \textbf{equa SS o de estrutura de Cartan} (ou equa SS o de Maurer--Cartan para fibrados principais).

A curvatura satisfaz a \textbf{identidade de Bianchi}:
\begin{equation}
 d\Omega + \omega\wedge\Omega - \Omega\wedge\omega = 0,
\end{equation}
que na GAT corresponde condi SS o de aus ancia de arbitragem (NFLVR): a curvatura do fibrado de pre SSos deve ser uma 2-forma fechada com respeito derivada covariante, garantindo que a carteira de ativos n o admite oportunidades de lucro sem risco. Este v -nculo geom ctrico entre curvatura e equil -brio de mercado constitui o cerne da Abordagem de Apre SSamento por Transporte Geom ctrico.

Para um tratamento completo da teoria de fibrados e suas aplica SS ues f -sica e finan SSas, remetemos o leitor a \citeonline{nakahara2003geometry}, \citeonline{kobayashi1963foundations} e \citeonline{bleecker1981gauge}.

\endinput


